% from ega4-18-fr.tex
\oldpage[IV-4]{110}

de $B$, $\mathfrak q_0$ est l'image d'un idéal premier $\mathfrak q$ de $B$ et $\mathfrak q$ l'image réciproque de $\mathfrak q_0$. Pour tout $i$, soit $v_i\in B$ un élément dont $u_i$ est l'image dans $B_0$, et soit $\mathfrak J$ l'idéal de $B$ engendré par les $v_i$, de sorte que $\mathfrak J_0$ est l'image de $\mathfrak J$ dans $B_0$. La proposition sera établie en prenant pour $U$ un voisinage ouvert du point de $\Spec(B/\mathfrak J)$ correspondant à l'idéal premier $\mathfrak p=\mathfrak q/\mathfrak J$, pourvu que l'on prouve que $B_{\mathfrak q}/\mathfrak J_{\mathfrak q}$ est une $A$-algèbre formellement lisse pour les topologies discrètes. Or, cela résulte du critère jacobien, car les images des $v_i$ dans $\mathfrak J_{\mathfrak q}/\mathfrak J_{\mathfrak q}^2$ engendrent ce $B_{\mathfrak q}$-module, et il résulte de (18.1.1.1) que l'on a $\det(\partial v_i/\partial T_{j_h})\notin\mathfrak q$.

\begin{theorem}[18.1.2]
\label{IV.18.1.2}
Soient $S$ un préschéma, $S_0$ un sous-préschéma fermé de $S$ dont l'espace sous-jacent est identique à celui de $S$. Alors le foncteur
\[
  X\longmapsto X\times_S S_0
\]
de la catégorie des $S$-préschémas étales sur $S$, dans la catégorie des $S_0$-préschémas étales sur $S_0$, est une équivalence de catégories.
\end{theorem}

Montrons d'abord que ce foncteur est pleinement fidèle. Soient $X$, $Y$ deux $S$-préschémas étales sur $S$, et posons $X_0=X\times_S S_0$, $Y_0=Y\times_S S_0$. Si $Z=X\times_S Y$, l'ensemble $\Hom_S(X,Y)$ est en correspondance biunivoque canonique avec l'ensemble des $X$-sections $\Gamma(Z/X)$, et de même $\Hom_{S_0}(X_0,Y_0)$ est en correspondance biunivoque canonique avec $\Gamma(Z_0/X_0)$, où $Z_0=Z\times_S S_0=X_0\times_{S_0}Y_0$. Or $Z$ est étale sur $X$, $Z_0$ étale sur $X_0$, et $X_0$ (resp. $Z_0$) est un sous-préschéma fermé de $X$ (resp. $Z$) ayant même espace sous-jacent. Les parties ouvertes de $Z$ telles que la restriction du morphisme $Z\to X$ soit surjective et radicielle sont donc les mêmes que les parties de $Z_0$ ayant les propriétés correspondantes et notre assertion découle par suite de (17.9.3).

Pour achever la démonstration, il suffit de voir que pour tout $S_0$-préschéma $X_0$ étale sur $S_0$, il existe un $S$-préschéma $X$ étale sur $S$ et un $S_0$-isomorphisme $X_0\xrightarrow{\sim}X\times_S S_0$. En vertu de la prop. (18.1.1), il existe un recouvrement ouvert $(U_\alpha)$ de $X_0$ et pour tout $\alpha$, un $S$-préschéma $V_\alpha$ qui est étale sur $S$, et enfin un $S_0$-isomorphisme $\theta_\alpha:U_\alpha\xrightarrow{\sim}V_\alpha\times_S S_0$. En outre, en raison de la première partie de la démonstration, il existe un unique $S_0$-isomorphisme $\varphi_{\alpha\beta}$ de $\theta_\alpha(U_\alpha\cap U_\beta)$ sur $\theta_\beta(U_\alpha\cap U_\beta)$, correspondant à l'automorphisme identique de $U_\alpha\cap U_\beta$, et il est immédiat, pour la même raison, que ces isomorphismes vérifient la condition de recollement ($0_{\mathrm I}$, 4.1.7). Il y a par suite un $S$-préschéma $X$ tel que les $V_\alpha$ s'identifient canoniquement à des sous-préschémas induits sur des ouverts de $X$, les $\theta_\alpha$ s'identifiant à des $S_0$-isomorphismes qui coïncident dans les intersections $U_\alpha\cap U_\beta$, et définissent donc un $S_0$-isomorphisme $X_0\xrightarrow{\sim}X\times_S S_0$. Il est clair que $X$ est étale sur $S$ (17.3.2), ce qui achève la démonstration.

\begin{corollary}[18.1.3]
\label{IV.18.1.3}
Soient $S$ un préschéma, $X$ un $S$-préschéma étale sur $S$, $S'$ un $S$-préschéma, $S'_0$ un sous-préschéma fermé de $S'$ ayant même espace sous-jacent. Alors l'application canonique $X(S')_S\to X(S'_0)_S$ (I, 3.4.3) est bijective.
\end{corollary}

Posons en effet $X'=X\times_S S'$, $X'_0=X\times_S S'_0$, de sorte que $X(S')_S=\Gamma(X'/S')$ et $X(S'_0)_S=\Gamma(X'_0/S'_0)$; le corollaire résulte du fait que le foncteur défini dans (18.1.2) (avec $S$ et $S_0$ remplacés par $S'$ et $S'_0$ respectivement) est pleinement fidèle (ou directement de (17.9.3)).

\subsection{Revêtements étales.}
\EGAReaderTarget{ega4.c04.subsec.18-2}{subsection:IV.18.2}

